Coverage-Guided Grey-box fuzzing is one of the most effective techniques for automatic vulnerability discovery in software
    \cite{zafl,Boehme2016},
    due to its ability to generate a wide range of inputs based on the feedback obtained through Code Coverage information,
    which guides the fuzzer towards unexplored regions of code and therefore helps it find hidden vulnerabilities.
    In fact, \textit{Black-box} fuzzing lacks a precise mutation strategy based on execution monitoring, while
    \textit{White-box} fuzzing incurs high overheads due to Symbolic or Concolic execution.
    Furthermore, the latter approach also requires the availability of the source code, which is not always available.
    So \textit{Grey-box} fuzzing seems to be the most promising trade-off between precision of input mutation and overhead.
    However, \textit{Grey-box} fuzzing also has its own drawbacks.
    As a matter of fact, the effectiveness of this approach strongly depends on the availability of the source code.
    In fact, both \gls{SBI} and \gls{DBI} have many limitations.
    The main limitation of \gls{DBI} is the huge amount of overhead (higher than 600\% in the case of AFL++ in QEMU
    mode \cite{zafl}), which makes the exploration of a substantial part
    of the valid input space almost impossible.
    Instead, the main difficulty of \gls{SBI} arises in modifying already compiled binaries.
    There are several problems regarding Binary Rewriting, such as Indirect Control Transfers (e.g. \texttt{call rax}),
    shifts in binary sections, or instructions that require rewriting the instructions that reference these sections, and many others,
    as seen in the previous \hyperref[sec:biniary_rewriting]{Section 1.2}.\\
    For all these reasons, the preferred technique nowadays is \textit{Compile-time Edge Coverage instrumentation} with fuzzers such as
    AFL++, LibAFL \cite{libafl}, or libFuzzer \cite{libfuzzer}.
    However, in practice, for reasons of industrial secrecy (i.e. proprietary software) or due to precompiled libraries, applications are
    often closed source, and so Compile-time Edge Coverage instrumentation becomes unusable.\\
    In these scenarios, binary-only fuzzing is the only viable approach.
    Although \gls{SBI} presents many challenges, it nonetheless proves to be the most fruitful area,
    given its low overheads \cite{zafl}.
    In this thesis I will address the problem
    of statically instrumenting a binary with Edge Coverage instructions in a way that makes the instrumented binary compatible with AFL++,
    which has been chosen as the fuzzer due to its popularity and because it represents the current state of the art, implementing
    most of the recently discovered techniques demonstrated by Andrea Fioraldi et al. \cite{aflplusplus}.\\
    The instrumentation must add as little overhead as possible with respect to compile-time overheads and must
    preserve the correctness of the binary. This is important in fuzzing, because lower overhead means more executions per second,
    so a larger portion of the valid input space can be explored.